\documentclass{article}

\usepackage[a4paper, left=1.5cm, right=1.5cm, top=2cm, bottom=2cm]{geometry}


\usepackage{hyperref}
\usepackage{fancyhdr}
\usepackage{ccicons}
\usepackage{caption}
\usepackage{subcaption}
\usepackage{algorithm}
\usepackage{algorithmic}

% Configuration des en-têtes et pieds de page
\pagestyle{fancy}
\fancyhf{} % reset tout

\fancyhead[L]{DL3 Math-Info}
\fancyhead[C]{CP UFR MATH}
\fancyhead[R]{2026-2027}

\fancyfoot[L]{Laurent CAI\\}
\fancyfoot[R]{\thepage}

\begin{document}

  \begin{center}
    \vspace{2em}
    {\huge\textbf{Commission pédagogique}}
    \vspace{1em}
  \end{center}

Est ressorti de cette séance de préparation les points suivants.

\section{But et propos liminaire}

\begin{itemize}
    \item Définir les aménagements de Licence 1 au Master 1 si les examens ont lieu en présentiel.
    \item Il s'agit de fournir les lignes générales qui seront ensuites déclinées dans chaque UE.
\end{itemize}

\noindent Rappel fonctionnement des commissions : \textbf{CP $\to$ CUFR $\to$ CFO/CFVU}. C'est les CFO/CFVU qui valident les aménagements. La plupart des changements proposer en réunion des MCCC sont a priori décidables par la CP.\\
\textbf{Attention : } Pour tous ceux qui pensent pouvoir bénéficier du relai handicap, manifestez-vous \textbf{maintenant} auprès de nous (représentants).\\
\textbf{Dispense et étudiants étrangers : } Celle dont il est question ne valide pas les crédits ECTS d'office (il faut valider le semestre 6 pour récupérer les ECTS). Il permet aux étudiants étrangers de ne pas redoubler et de renouveller leur titre de séjour. Idem \textit{a priori} pour les étudiants boursiés.\\\\
Il est possible pour les étudiants étrangers concernés de demander une lettre à l'UFR à donner à la préfeture. Les enseignants précisent qu'ils le feront, mais que le préfet à tendance à ne pas écouter.


\section{Cas des Maths-Info}

\textit{Il s'agit d'un résumé de ce que la CP a validé de fait pour les DL3 Maths-Info. Cela n'a pas été discuté lors de la CP.}

\begin{itemize}
    \item Deux évaluations en présentiel pour le semestre : \textbf{l'examen de janvier, et un contrôle continu} (a priori organisé plus tard) ;
    \item Le contenu des évaluations est indiqué par les enseignants sur la page Moodle, en particulier :
    \begin{itemize}
        \item \textbf{Calcul différentiel :} L'examen et le CC porteront sur le programme de ce qui a été fait en présentiel. Éventuellement, une part de différentiabilité sera en \textbf{bonus} pour ceux qui ont suivi le cours ;
        \item \textbf{Intégration et probabilités :} L'examen portera sur tout ce qui a été fait, distanciel y compris, avec des aménagements de sorte que le présentiel soit suffisant pour obtenir une bonne note ;
        \item \textbf{Groupes et actions :} L'examen portera sur tout ce qui a été fait, distanciel y compris, avec des aménagements de sorte que le présentiel soit suffisant pour obtenir une bonne note.
        \item \textit{les fiches contenant le programme exact de chaque U.E. sont dores et déjà disponibles sur Moodle.}
    \end{itemize}
    \item La note de première session du semestre sera : $NF = \frac{CC + E}{2}$, avec $CC$ la note du contrôle continu et $E$ la note de l'examen de janvier ;
    \item Les copies seront corrigées de manière anonyme, et décachetées par l'enseignant devant un étudiant ;
    \item Les étudiants peuvent demander des dispenses dans chaque matière, le jury "étudiera les demandes avec une grande bienveillance".
\end{itemize}

\section{Aménagements généraux}
\begin{itemize}
    \item Toute évaluation doit être prévenue une semaine à l'avance au moins \textit{(empêchant de programmer des évaluations sur une semaine de réouverture)} ;
    \item Un programme précis et détaillé des évaluations doit être fourni dans chaque UE, il peut être allégé ;
    \item Ce programme doit expliciter les parties initialements prévues qui sont supprimées ;
    \item La partie des évaluations qui porte sur le programme vu en présentiel doit permettre d'obtenir une très bonne note;
    \item Le jury est souverain, et il étudiera avec bienveillance les demandes de dispense ;
    \item \textbf{Si pas d'examen en présentiel, ils vont commencer à réfléchir}, ils espèrent aboutir avant les vacances de Noël à une discussion. La délocalisation est sur la table.
\end{itemize}

\section{Accord sur le principe sans engagement}

\begin{itemize}
    \item Assurer l'anonymat des copies pendant des copies examens ; 
    \item Le deuxième semestre rependra les bases qui n'auront pas été vues au Semestre 5 \textit{(il y aura plus de rappels)} ;
    \item Il sera possible de demander le redoublement en Maths-Info ; les demandes seront évaluées de manière favorable d'après Monsieur Malod.
\end{itemize}

\section{Motion votée}
\textbf{Les étudiants invités n'avaient pas le droit de vote. La motion suivant a été adoptée à l'unanimité par les enseignants votants présents.}\\

\noindent La Commission Pédagogique de l'UFR de Mathématiques réunie le mardi 24 novembre 2026 a voté l'organisation générale des épreuves de fin de semestre de ses formations. Les dispositions suivantes sont adoptées concernant les MCCC en Licence et Master 1:

\begin{enumerate}
    \item Les UE soumises au Contrôle Continu Intégral comporteront toujours au moins deux évaluations, comptant chacune pour au plus 50\% de la note finale. Au cas par cas selon les UE et les filières et sous réserve d'acceptation des conseils, le nombre d'épreuves de contrôle continu pourra être réduit, la nature de certaines épreuves pourra être modifiée, et les épreuves de seconde chance pourront être passées par tous les étudiants. Les nouvelles modalités de contrôle des connaissance devront être transmises aux étudiants au plus vite UE par UE et filière par filière.

    \item La Commission Pédagogique exprime son accord pour que les filières le souhaitant passent du CET au CT, sous réserve d'acceptation de la CFO. Les filières en CCI n'ayant pas pu avoir de Contrôle Continu pourront demander au CFO de passer en CT.

    \item La Commission Pédagogique demande aux enseignants de corriger des copies anonymisées.

    \item Toute évaluation, de quelque nature qu'elle soit, doit avoir été programmée et communiquée aux étudiants au moins une semaine à l'avance.

    \item Les épreuves terminales des UE de Licence hors MIASHS ayant commencé avant les blocages comporteront une partie sur les enseignements dispensés avant les blocages (communs à tous les groupes pour les cours/TD de L1) ou en présentiel en cas de retour à la normale à temps. Cette partie devra permettre d'obtenir une très bonne note. Le cas spécifique de l'évaluation des UE sans enseignement présentiel avant le blocage sera pris en compte dans chaque parcours. En Licence MIASHS et en M1, les enseignants suivront les résultats des consultations avec les étudiants filière par filière.

    \item Le programme précis et détaillé des évaluations de chaque UE devra être communiqué au plus tôt aux étudiants, UE par UE et filière par filière. Il devra être clairement explicité quels contenus initialement prévus dans l'UE seront exclus du programme des évaluations.

    \item Les jurys de chaque filière sont souverains et cette commission ne saurait s'y substituer. Elle recommande cependant que les circonstances exceptionnelles dans lesquelles se déroule ce semestre soient prises en compte avec bienveillance.
\end{enumerate}

Pour les parcours de M2, l'organisation des examens devra être discutée en accord avec les autres établissements cohabilités.


\end{document}
\end{document}